\subsection{Rank-one sources with two bad Weierstrass transfer fibres}

\paragraph{Finding A43 (the rank-one and Weierstrass qualifications do not repair the fixed-cover detector).}
Let $E/\mathbf F_2$ be the ordinary curve of A13 and A38, and
base change it to $k=\mathbf F_{128}$. In the inverse family
of A40 there are at least $95$ parameters $t\in E(k)$ for
which the smooth cyclic degree-five cover
$\pi_t:D_t\to C_t$ has the following properties:
\[
 \begin{gathered}
       g(C_t)=2,\qquad f(C_t)=1,\\
 W_t^{(1)},W_t^{(2)}\in C_t(k)\text{ are both Weierstrass points,}\\
 \text{both fibres are split in }D_t,\\
 \operatorname{ord}([K_E-2R])=4159,\\
 \text{for every lift of either point and every reflection image }R.
 \end{gathered}
 \tag{RWI.1}
\]
Since $\operatorname{ord}_{4159}(2)=2079$ is odd, none of
these images satisfies the sufficient hypothesis of A29.
Their exact canonical orders persist after every extension
of constants. Thus changing the lift or reflection of this
fixed cover does not force the canonical detector even for
a rank-one source and either prescribed Weierstrass point.
The assertion does not quantify over other cyclic covers
of that source point.

\begin{proof}
\emph{1. The special Frobenius and automorphism identities.}
Write $F$ for the geometric-point map obtained by squaring
the coordinates on $E$, and write $\pi$ for the corresponding
Frobenius endomorphism of $J=J_E$. Put $V=\pi^\dagger$.
The checked A33 supplement supplies the integral,
Rosati-symmetric endomorphism $T$ with
\[
 \pi V=2,\qquad \pi+V=2T+1,\qquad T^2+T=1.
 \tag{RWI.2}
\]
All these endomorphisms are defined over $\mathbf F_2$ and
commute. The element
\[
 u=T-\pi
 \quad\text{satisfies}\quad
 u^\dagger u=1,\qquad u^2+u+1=0.
 \tag{RWI.3}
\]
Indeed $u^\dagger=T-V$, and both identities follow by
substitution in (RWI.2). It is an integral polarized
automorphism of exact order three. The genus-two Torelli
uniqueness statement and descent give an automorphism
$\rho:E\to E$ over $\mathbf F_2$ inducing $u$; its order
is three by faithfulness of the curve action on its
Jacobian \cite[Theorem~12.1 and Corollary~12.2,
pp.~38--39]{MilneJV}.

Use the actual Raynaud theta characteristic $B$ on $E$;
the Frobenius twist is identified with $E$ over
$\mathbf F_2$. It is defined over $\mathbf F_2$, is preserved
by $\rho$, and is non-effective because $E$ is ordinary.
Put $z_P=[P-B]$ and $\Theta=\{z_P:P\in E\}$.
Thus $0\notin\Theta$, $z_{\iota P}=-z_P$,
$\pi z_P=z_{F(P)}$, and $u z_P=z_{\rho(P)}$.
The three Weierstrass classes are exactly the three
nonzero geometric points of $J[2]$ on $\Theta$, as in A35.
They are all $\mathbf F_2$-rational. Since $T$ is an
automorphism, the intersection number from (ID.4) implies
\[
             \Theta\cap T\Theta
                =J[2](\bar k)\setminus\{0\}
             \quad\text{as a set of three points}.
 \tag{RWI.4}
\]
Each already belongs to both divisors, whose intersection
number is three.

Changing the theta characteristic does not change the
residual family. Intrinsically, for $P\ne t$ its fibre is
$Q+R\sim K_E+T([P-t])$, as in (ID.7); equivalently, if
two theta characteristics differ by $a\in J[2]$, the two
maps and theta divisors in (ID.5)--(ID.6) are translated
by the same $a$. Hence the family used here is the one
counted in A41.

Differentiating (RWI.2) gives $dV=1$, since $d\pi=0$
and $d[2]=0$. In particular $V$ is etale. Its degree is
four and its kernel is all four geometric points of
$J[2]$: $\pi$ acts identically on those points and
$V\pi=[2]$. Also $d\rho=du=dT$.
Let $\gamma:E\to\mathbf P(T_0J)$ be the canonical Gauss
map. In an $\mathbf F_2$ coordinate its three branch
values are exactly $\mathbf P^1(\mathbf F_2)$, because
the branch values of the displayed A13 model are
$0,1,\infty$. The projective action of $dT=d\rho$
permutes those three values.

\emph{2. The normalized double-point locus.}
The image
\[
                         \Gamma=V\Theta\subset J
 \tag{RWI.5}
\]
has normalization $E$. Indeed $V$ is etale, and the only
identifications on $\Theta$ by its nonzero kernel points
are pairs of Weierstrass classes, by the theta-intersection
calculation in A35. All three Weierstrass classes map to
zero. Their tangent lines are distinct and $dV=1$, so
$\Gamma$ has three smooth transverse branches at zero
and is smooth away from zero. Numerically $\Gamma\equiv2\Theta$:
$\pi_*\Theta=2\Theta$ and $[2]_*\Theta=4\Theta$, while
$[2]=V\pi$ and $V|_\Theta$ is birational. The pushforward
identity follows from the multiplication pullback and projection
formulas \cite[Corollary~6.6 and Lemma~8.3]{MilneAV}.
Thus $\Gamma^2=8$, whereas $(T\Theta)^2=2$, so $\Gamma$
cannot coincide with a translate of $T\Theta$.

The condition that a degree-two effective divisor is a
double point is the image of $Q\mapsto 2z_Q$.
Because $[2]=V\pi$ and $F:E\to E$ is a radicial degree-two
map, its reduced image is precisely $\Gamma$.
For non-Weierstrass $t$ put
$H_t=T(\Theta-z_t)$. The total intersection number is six.
For example the Rosati trace formula gives
\[
 \begin{aligned}
 H_t\cdot\Gamma
 &=\Theta\cdot V^*(T\Theta)\\
 &=\tfrac12\operatorname{Tr}
       \bigl((T^{-1}V)^\dagger(T^{-1}V)\bigr)\\
 &=\tfrac12\operatorname{Tr}\bigl(2(T+2)\bigr)=6.
 \end{aligned}
 \tag{RWI.6}
\]
Here $T^{-2}=T+2$ and the rational Tate trace of $T$ is
$-2$. At zero the curve $H_t$ has tangent
$dT\gamma(t)$, which is not one of the three branch
lines of $\Gamma$, since $t$ is not Weierstrass.
Thus zero contributes exactly three to (RWI.6).

We determine the remaining intersection multiplicities.
On $E\times E$ use the normalized map
\[
                    \Psi(P,Q)=Tz_P-Vz_Q.
 \tag{RWI.7}
\]
Away from zero, a tangency of $H_t$ and $\Gamma$ is
characterized by $\gamma(Q)=dT\gamma(P)$, since
$dV=1$. Its two critical graphs are
$Q=\rho(P)$ and $Q=\iota\rho(P)$.
The identities in (RWI.2)--(RWI.3) give
\[
 \begin{aligned}
 T-Vu&=-u^2\pi^2,\\
 T+Vu&=-T\pi.
 \end{aligned}
 \tag{RWI.8}
\]
On the first graph, $\Psi(P,Q)=-u^2\pi^2z_P$ lies in
$\Theta$. It cannot equal $Tz_t$ for non-Weierstrass $t$,
by (RWI.4). On the second graph it equals $-T\pi z_P$.
Consequently there is exactly one possible nonzero
point of tangency, specified by
\[
                   t=\iota F(P),\qquad Q=\iota\rho(P).
 \tag{RWI.9}
\]
These points are non-Weierstrass and hence $Vz_Q\ne0$.

The tangency has multiplicity exactly two, not three.
At the point in (RWI.9), the differential of $\Psi$ has
one-dimensional image $dT\gamma(P)$. The tangent of
$T\Theta$ at $Tz_t$ is $dT\gamma(F(P))$.
These lines differ: equality would make $\gamma(P)$
fixed by Frobenius on $\mathbf P^1$, hence a branch
value and $P$ Weierstrass. The inverse image
$\Psi^{-1}(T\Theta)$ is therefore a smooth reduced
curve locally at this point. It contains the smooth
second critical graph, so locally equals that graph.
Its map to $T\Theta$ is, by (RWI.9), the radicial
Frobenius map of degree two. Thus its local fibre has
length exactly two. Since $V\Theta$ is smooth here
and its normalization is an isomorphism locally, this
is exactly the local intersection length of $H_t$
with $\Gamma$.

After the origin contribution three and this contribution
two, (RWI.6) leaves exactly one further simple nonzero
intersection. For a smooth $D_t$, these two nonzero
intersections are exactly the ramification points of
$v_1:D_t\to E$. Off $P=t$, this follows directly from
the unique effective pair in (ID.7): its two points
coincide precisely on the double-point locus. Over
$P=t$, the residual fibre is characterized by
$\gamma(Q)=dT\gamma(t)$; it consists of two distinct
points because this value is not a branch value.
There is therefore no further ramification there.

\emph{3. The source has two-rank one.}
The degree-two Galois cover $v_1$ has exactly two
geometric ramification points. Deuring--Shafarevich gives
\[
                 f(D_t)-1=2(f(E)-1)+2=4,
                      \qquad f(D_t)=5
 \tag{RWI.10}
\]
\cite[Corollary~1.8, pp.~35--36]{Crew1984}.
The isogeny $J_{D_t}\sim J_E^2\times J_{C_t}$ from
A40 and invariance and additivity of two-rank under
isogeny give $5=4+f(C_t)$, so $f(C_t)=1$.
The excluded vertical-branch parameter set in A41
contains the three Weierstrass parameters, since
$dT$ permutes their Gauss values. Its entire excluded
set still has at most $34$ points. Thus at least $95$
parameters in $E(\mathbf F_{128})$ are both smooth and
non-Weierstrass, and all have this rank-one source.

\emph{4. A reflection-fixed lift with all five images bad.}
Fix such a parameter $t$, and take
$P=\iota F^{-1}(t)\in E(k)$. Set
\[
 \begin{gathered}
 x=z_P,\qquad z_t=-\pi x,\qquad
 y=-u\pi^{-1}x=z_{\iota\rho F^{-1}(P)},\\
 x'=u^2\pi^{-1}x=z_{\rho^2F^{-1}(P)}.
 \end{gathered}
 \tag{RWI.11}
\]
The inverse $\pi^{-1}$ here means inverse Frobenius
on geometric points; on $k$-points it is $\pi^6$.
All three displayed classes are points of $\Theta(k)$
and correspond to non-Weierstrass points. They are
pairwise distinct. For example $x=y$ would imply
$Tx=0$, $x=x'$ would imply $(T+1)x=0$, and $x'=y$
would imply $x=0$. These are impossible because $T$
and $T+1$ are automorphisms and $0\notin\Theta$.

None of their points equals the auxiliary parameter
$t$. For $x=z_t$, (RWI.9) and $Vz_Q\ne0$ already
exclude the equality. If $x'=z_t$, then
$\pi^2x=-u^2x$; applying the seventh power and using
$\pi^7x=x$ gives $x=-u^2x$, hence $ux=0$.
If $y=z_t$, then $\pi^2x=ux$; the seventh power
implies $ux=x$ and then $\pi^2x=x$. The corresponding
point of $E$ is in $E(\mathbf F_4)$, whose only points
are the three Weierstrass points by A13. Again this
contradicts the choice of $P$.

For two theta classes $a,b$, abbreviate
$F_t(a,b)=Ta-b-Tz_t$. Direct use of (RWI.2)--(RWI.3)
gives
\[
 \begin{gathered}
 F_t(x,y)=y,\qquad F_t(x',y)=x',\qquad
 F_t(y,x')=x. 
 \end{gathered}
 \tag{RWI.12}
\]
All first coordinates are different from $t$, so the
three pairs in (RWI.12) belong to the residual curve,
not just to its removed vertical component. Let
$a_0=(x,y)\in D_t(k)$, identifying theta classes with
their unique points of $E$. The first equality makes
$a_0$ fixed by $\tau_1$, the involution of $v_1$.
The other two, and uniqueness of the other point in
each degree-two fibre, give the five distinct points
\[
 \begin{gathered}
 a_0=(x,y),\quad \sigma a_0=(x',x'),\quad
 \sigma^2a_0=(y,x),\\
 \sigma^3a_0=(y,x'),\quad\sigma^4a_0=(x',y),
             \qquad \sigma=\tau_1\tau_2.
 \end{gathered}
 \tag{RWI.13}
\]
For instance $\tau_2(x,y)=(x',y)$ and
$\tau_1(x',y)=(x',x')$; the next step uses $F_t(y,x')=x$.
The last two orbit points also follow directly from
$\tau_1\sigma^j a_0=\sigma^{-j}a_0$, since
$\tau_1a_0=a_0$, by swapping the first-projection fibre.

Put $W_t^{(1)}=\pi_t(a_0)$. It is rational and its entire
cyclic degree-five fibre is the rational orbit
(RWI.13), so it is split. Since $\tau_1a_0=a_0$
and $\tau_1$ descends to the hyperelliptic involution
of $C_t$, the point $W_t^{(1)}$ is Weierstrass.
The five reflection images are the first coordinates
of (RWI.13), namely $x,x',y,y,x'$.

\emph{5. The other Weierstrass fibre has the same property.}
Put $a=z_t$ and define three further theta points
\[
       x_2=-\pi a,\qquad y_2=u\pi^{-1}a,
                       \qquad x'_2=u^2\pi^{-1}a.
 \tag{RWI.14}
\]
They are respectively the classes of
$\iota F(t)$, $\rho F^{-1}(t)$ and
$\rho^2F^{-1}(t)$, so all are non-Weierstrass and
$k$-rational. They are pairwise distinct. Equality
of $x_2$ with $y_2$ or $x'_2$ would give
$\pi^2a=-u^j a$ for $j=1$ or $2$. Raising to the
seventh power shows $\pi^2a=a$, putting $t$ in
$E(\mathbf F_4)$, which is impossible. Equality
$y_2=x'_2$ would give $ua=a$, hence $3a=0$ by
(RWI.3). But $\#J(k)=4\cdot4159$ by A38, so this
would force $a=0$, again impossible.

None of these three points is $t$. Equality $x_2=a$
implies $\pi^2a=a$ and is excluded as above.
Equality $y_2=a$ implies $\pi a=ua$; taking seventh
powers gives $ua=a$, excluded by the same absence
of rational three-torsion. Finally $x'_2=a$
implies $(\pi-u^2)a=(T+1)a=0$, contradicting
invertibility of $T+1$.

The same ring calculations give
\[
 F_t(x_2,y_2)=y_2,\qquad
 F_t(x'_2,y_2)=x'_2,\qquad
 F_t(y_2,x'_2)=x_2.
 \tag{RWI.15}
\]
Consequently $a_2=(x_2,y_2)$ is another rational
$\tau_1$-fixed point of $D_t$, and its orbit is
exactly (RWI.13) with every variable replaced by
its subscript-two version. Its source point
$W_t^{(2)}=\pi_t(a_2)$ is therefore rational,
Weierstrass and split, with all five reflection
images non-Weierstrass.

The two fixed lifts are distinct: their first
coordinates are $\iota F^{-1}(t)$ and $\iota F(t)$,
whose equality would again put $t$ in $E(\mathbf F_4)$.
A reflection has at most one fixed point in a free
$C_5$ fibre. Indeed if it fixes $b$ and $\sigma^j b$,
then inversion of the deck group gives
$\sigma^{2j}b=b$, forcing $j=0$. Hence
$W_t^{(1)}\ne W_t^{(2)}$.
The hyperelliptic cover of a rank-one genus-two curve
has exactly two geometric branch points, by the same
Deuring--Shafarevich formula. These are therefore both
Weierstrass points of $C_t$.

All images in both orbits are non-Weierstrass points
of $E(k)$, so A38 gives exact canonical order $4159$
for each. Changing the lift only permutes its orbit,
and changing the reflection runs through the same
five images. An isomorphism of a quotient with $E$
preserves canonical-class order. Finally constant
extension preserves the order of each fixed geometric
class. This proves (RWI.1) and the stated detector failure.
\end{proof}
